\chapter{Getting Started}

\api{crystal-qt6} is a focused Qt6 Widgets binding for Crystal. It favors an explicit, tested desktop surface over automatic binding generation, so each exposed class has a Crystal wrapper, specs, and at least one practical usage path.

\section{Why Crystal?}

Crystal is attractive for desktop application work because it gives you a compact, expressive syntax while still compiling to native code. A Crystal program can feel close to Ruby while keeping static types, ahead-of-time compilation, and direct C interop. That combination matters for GUI code: you want application logic to stay readable, but you also want the process to start quickly, call native libraries directly, and package like a compiled program.

The type system is also a useful design partner. Widget state, model data, callback signatures, and value objects such as colors, rectangles, fonts, and indexes all benefit from APIs that say what they accept and return. The result is a programming style that remains pleasant for small examples but scales better as the interface grows.

\section{Why Qt?}

Qt is one of the most mature cross-platform GUI toolkits available. Its Widgets stack includes native windows, layouts, menus, actions, dialogs, item views, text editing, painting, images, SVG, PDF output, timers, clipboard access, drag and drop, and much more. These are not separate pieces that merely happen to coexist; they are designed to work together inside one event loop and object model.

That depth is the reason \api{crystal-qt6} is useful. Instead of building a small custom UI layer and then discovering that file dialogs, model/view tables, undo stacks, text editing, or image rendering are missing, the \api{crystal-qt6} library grows around Qt's existing desktop vocabulary. The goal is not to hide Qt. The goal is to make Qt feel natural from Crystal.

\section{What This Binding Is Not Yet}

\api{crystal-qt6} is a growing Qt Widgets binding, not a complete replacement for PySide6, PyQt, or the full C++ Qt API. It intentionally exposes a focused desktop surface first: widgets, layouts, application shells, dialogs, events, painting, images, undo stacks, model/view, clipboard data, and maintained examples.

That focus is useful, but it also means you should expect gaps. Some Qt modules are not wrapped. Some advanced overloads are omitted until there is a clear Crystal-facing use case. Some APIs are shaped around explicit ownership and typed values rather than exact source-level compatibility with another Qt binding.

When you are deciding whether to build on this shard, treat the examples and specs as the source of truth for the current surface. If a class or method is documented in this guide, it should have a maintained usage path. If you need a missing Qt feature, the usual path is to add a small native bridge function, wrap it in \api{src/qt6}, add specs, and then update the relevant example or chapter.

\section{Requirements}

You need Crystal, Qt6 development packages visible through \api{pkg-config}, and a C++ compiler compatible with the platform Qt build. The repository build currently expects Qt Widgets, Qt Svg, and Qt Svg Widgets.

\section{Installing On macOS}

The most direct macOS setup uses Homebrew. Install Apple's command line tools first if they are not already present:

\begin{lstlisting}[style=crystal]
xcode-select --install
\end{lstlisting}

Then install Crystal, Qt6, and \api{pkg-config}:

\begin{lstlisting}[style=crystal]
brew update
brew install crystal qt pkg-config
\end{lstlisting}

Homebrew's Qt formula is also known as \api{qt6} or \api{qt@6}, but \api{qt} is the usual package name. On most Homebrew installations, \api{pkg-config} can find Qt after installation. If it cannot, add Qt's package configuration directory to \api{PKG_CONFIG_PATH}:

\begin{lstlisting}[style=crystal]
export PKG_CONFIG_PATH="$(brew --prefix qt)/lib/pkgconfig:${PKG_CONFIG_PATH}"
\end{lstlisting}

Put that export in your shell startup file if the verification step below only works after setting it manually.

\section{Installing On Linux}

Crystal provides an official installer for DEB- and RPM-based Linux distributions. It configures the Crystal package repository and installs the \api{crystal} package:

\begin{lstlisting}[style=crystal]
curl -fsSL https://crystal-lang.org/install.sh | sudo bash
\end{lstlisting}

After Crystal is installed, install a compiler, \api{make}, \api{pkg-config}, and the Qt6 development packages that provide \api{Qt6Widgets}, \api{Qt6Svg}, and \api{Qt6SvgWidgets}.

On Debian or Ubuntu:

\begin{lstlisting}[style=crystal]
sudo apt update
sudo apt install -y build-essential pkg-config qt6-base-dev libqt6svg6-dev
\end{lstlisting}

On Fedora, Red Hat, or compatible \api{dnf}-based systems:

\begin{lstlisting}[style=crystal]
sudo dnf install -y gcc-c++ make pkgconf-pkg-config qt6-qtbase-devel qt6-qtsvg-devel
\end{lstlisting}

On openSUSE or SUSE Linux Enterprise:

\begin{lstlisting}[style=crystal]
sudo zypper install -y gcc-c++ make pkgconf-pkg-config qt6-base-devel qt6-svg-devel
\end{lstlisting}

Distribution package names vary, especially for Qt modules. If your package manager uses different names, search for the Qt6 base and SVG development packages and verify that they install \api{pkg-config} metadata for all three Qt modules listed above.

\section{Verifying The Toolchain}

Check Crystal first:

\begin{lstlisting}[style=crystal]
crystal --version
\end{lstlisting}

Then confirm that \api{pkg-config} can resolve every Qt module required by the native bridge:

\begin{lstlisting}[style=crystal]
pkg-config --modversion Qt6Widgets Qt6Svg Qt6SvgWidgets
\end{lstlisting}

That command should print version numbers without errors. If it fails, Crystal may still be installed correctly, but the Qt development packages are not visible to \api{pkg-config}. Fix that before building the shard.

\section{First-Build Troubleshooting}

Most first-build failures come from the native Qt toolchain rather than Crystal source code. Start by rerunning the exact \api{pkg-config} command from the previous section.

\begin{description}
\item[\api{pkg-config} is missing.] Install \api{pkg-config} or your distribution's equivalent package. Fedora and related systems often call it \api{pkgconf-pkg-config}.
\item[\api{Qt6Widgets} is not found.] Install the Qt6 base development package. On Debian and Ubuntu this is usually \api{qt6-base-dev}; on Fedora it is \api{qt6-qtbase-devel}; on openSUSE it is \api{qt6-base-devel}.
\item[\api{Qt6Svg} or \api{Qt6SvgWidgets} is not found.] Install the Qt6 SVG development package. The binding uses SVG support for rendering examples and documentation-friendly widgets, so base Qt alone is not enough.
\item[Homebrew Qt.] If Qt is installed but not found, add Homebrew's Qt package configuration directories to \api{PKG_CONFIG_PATH}. Check both \api{lib/pkgconfig} and \api{share/pkgconfig}; the CI workflow uses the same idea.
\item[Crystal cannot write compiler cache files.] Point Crystal's cache at a writable directory:
\end{description}

\begin{lstlisting}[style=crystal]
export CRYSTAL_CACHE_DIR=/tmp/crystal-cache
\end{lstlisting}

After changing packages or environment variables, run \api{make native} once. That isolates native bridge errors from normal Crystal compile errors and gives shorter output while you are fixing the toolchain.

\section{A First Window}

Every Qt program has one application object and one event loop. The \api{Qt6.application} helper creates or returns that process-wide wrapper. Widgets are then created, shown, and kept responsive by entering \api{app.run}.

\begin{lstlisting}[style=crystal]
require "qt6"

app = Qt6.application

window = Qt6.window("Counter", 320, 140) do |widget|
  count = 0
  label = Qt6::Label.new("0")
  button = Qt6::PushButton.new("Increment")

  button.on_clicked do
    count += 1
    label.text = count.to_s
  end

  widget.vbox do |layout|
    layout << label
    layout << button
  end
end

window.show
app.run
\end{lstlisting}

The call to \api{window.show} makes the top-level widget visible. The call to \api{app.run} hands control to Qt's event loop. From that point on, user input, timers, painting, dialog results, and signal callbacks are delivered by Qt until the application exits.

\section{Lifecycle And Shutdown}

Crystal and Qt manage object lifetimes differently. Qt uses parent-owned objects: a window can own its menu bar, toolbars, child widgets, dialogs, timers, and layout objects. Crystal uses garbage collection for ordinary Crystal objects. The \api{crystal-qt6} library keeps these two models predictable by making native Qt cleanup deterministic instead of depending on Crystal finalizers.

For ordinary applications, create widgets on the GUI thread, give child widgets a parent when possible, show the top-level window, and let \api{app.run} return when the last window closes or when \api{app.quit} is called. The \api{crystal-qt6} library registers \api{Qt6.shutdown} for process exit so short examples do not need explicit teardown code.

\begin{lstlisting}[style=crystal]
quit = Qt6::Action.new("Quit", main)
quit.shortcut = "Ctrl+Q"
quit.on_triggered do
  app.quit
end
\end{lstlisting}

You may call \api{Qt6.shutdown} yourself at a clear application boundary, but most GUI programs should not need to. The important rule is to avoid treating Qt widgets like disposable background objects. Create and use them on the GUI thread, prefer parent ownership, and allow Qt to finish its event loop before process shutdown.

\section{Build Flow}

The native shim builds automatically before Crystal compiles the shard. The repository \api{Makefile} provides convenient shortcuts for common tasks:

\begin{lstlisting}[style=crystal]
# Build the native bridge directly.
make native

# Run the spec suite.
make spec

# Launch a maintained example.
make example-dialogs
\end{lstlisting}

Those targets are not special entry points into the library. They wrap ordinary Crystal commands, so direct commands work as well:

\begin{lstlisting}[style=crystal]
# Run the same spec suite without make.
crystal spec

# Compile or run an example directly.
crystal build examples/dialog_gallery.cr
crystal run examples/dialog_gallery.cr
\end{lstlisting}

For day-to-day work, use whichever form is clearer in context. The \api{make} targets are nice for discoverability and CI-style repetition; direct \api{crystal} commands are useful when you are experimenting with one file or passing compiler options.

\section{Running Without A Desktop}

GUI code still needs a Qt platform backend, even when it is running in automation. The repository CI uses two different approaches:

\begin{itemize}
\item Linux runs the specs under \api{xvfb-run} with \api{QT_QPA_PLATFORM=xcb}. This gives Qt an X11-like display without requiring a physical monitor.
\item macOS runs the specs with \api{QT_QPA_PLATFORM=offscreen}. This is useful for noninteractive widget tests, but it is not a substitute for visually checking examples that depend on native window behavior.
\end{itemize}

For local Linux automation, the CI command is a good starting point:

\begin{lstlisting}[style=crystal]
xvfb-run -a env QT_QPA_PLATFORM=xcb make spec
\end{lstlisting}

For local macOS smoke tests, \api{QT_QPA_PLATFORM=offscreen make spec} can be useful when you only need the specs to run. For screenshots, native dialogs, or layout inspection, run the examples in a normal desktop session.

\section{Project Layout}

The repository is small enough to navigate directly:

\begin{description}
\item[\api{src/qt6}] Crystal wrappers, enums, value objects, callback helpers, and public APIs.
\item[\api{ext/qt6cr}] The C++ native bridge that talks to Qt and exposes stable C-callable functions to Crystal.
\item[\api{examples}] Maintained runnable programs that exercise real usage paths. These are the best place to learn how pieces fit together.
\item[\api{spec}] End-to-end specs for wrappers, ownership behavior, callbacks, rendering, dialogs, models, and examples of supported edge cases.
\item[\api{docs/book}] This guide, its LaTeX source, generated PDF, and screenshot assets.
\item[\api{scripts}] Build and documentation helpers, including native bridge scripts and screenshot/figure generation.
\item[\api{.github/workflows}] CI setup for macOS and Linux.
\end{description}

If you are reading as an application developer, start with \api{examples} and the matching chapters. If you are reading as a contributor, pair a wrapper in \api{src/qt6} with its bridge functions under \api{ext/qt6cr} and its coverage in \api{spec}.

\section{How To Read This Guide}

The early chapters cover the object model and common widgets. Later chapters move into desktop-shell composition, dialogs, custom event handling, painting, image processing, undo stacks, and model/view patterns. Each topic includes a small API sketch and points to a maintained example when possible.

After this chapter, Chapter 2 is the most important conceptual stop: it explains how Qt ownership and Crystal wrappers cooperate. Chapter 3 then gives you the everyday widget and layout vocabulary. Chapter 5 is a good early payoff chapter because standard dialogs immediately make small applications feel complete. If you prefer to learn by following one thread from start to finish, skim Chapters 2 through 10 and then read Chapter 11 alongside \api{examples/editor_vertical_slice.cr}.
